
de façon que le
premier facteur permute @@M0@@ et le second @@M1@@~; @@M2@@ est plongé dans l'anneau des invariants en envoyant @@M3@@ sur la @@M4@@-ième fonction symétrique des @@M5@@~; @@M6@@ envoie la @@M7@@-ième fonction symétrique de @@M8@@ (resp. @@M9@@) sur le @@M10@@-ième coefficient de @@M11@@ (resp. @@M12@@). Alors~:
@@M13@@
Comme @@M14@@ est un homomorphisme de groupes, on est ramené au cas où @@M15@@, et dans ce cas l'additivité de @@M16@@ résulte de la définition des @@M17@@.

1.7. On peut généraliser les considérations précédentes au cas d'un produit de foncteurs @@M18@@~; le foncteur @@M19@@ est représentable par l'anneau
@@M20@@, @@M21@@, et @@M22@@, où @@M23@@ est de poids @@M24@@. On notera @@M25@@ la série @@M26@@~; l'élément @@M27@@ de @@M28@@ possède la propriété universelle analogue à celle de @@M29@@ (cf. 1.2).

On peut introduire les morphismes @@M30@@, @@M31@@ comme précédemment en refaisant la construction pour chaque famille @@M32@@ séparément. On obtient un résultat analogue au théorème 1.6 en cherchant les homomorphismes de foncteurs de @@M33@@ dans @@M34@@, @@M35@@-multilinéaires et compatibles aux filtrations par rapport à chacun des @@M36@@ arguments séparément : ils sont déterminés par leur valeur pour l'élément
